\chapter{The Schmidt decomposition and purification}\label{app:schmidt}

\section{Derivation}
Consider a pure state of a bipartite system, $\ket\psi_{AB}\in\Ha_A\otimes\Ha_B$. In arbitrary orthonormal bases $\{\ket i_A\}$ and $\{\ket\mu_B\}$ of the two subsystems, any such state can be written
\begin{equation}
  \ket\psi_{AB}=\sum_{i,\mu}\psi_{i\mu}\ket i_A\otimes\ket\mu_B .
\end{equation}
Grouping the terms with the same $\ket i_A$, we can also write
\begin{equation}
  \ket\psi_{AB}=\sum_i\ket i_A\otimes\ket{\tilde\imath}_B,\qquad\text{with}\qquad\ket{\tilde\imath}_B=\sum_\mu\psi_{i\mu}\ket\mu_B .
\end{equation}
In general the states $\ket{\tilde\imath}_B$ are neither orthogonal nor normalized. Let us show that they are orthogonal if the basis of $A$ is suitably chosen. We choose the basis $\{\ket i_A\}$ in which the reduced density matrix of $A$ is diagonal:
\begin{equation}
  \hat\rho_A=\sum_ip_i\proj i .
\end{equation}
On the other hand, by definition of the reduced density matrix,
\begin{equation}
  \hat\rho_A=\tr_B\proj\psi=\tr_B\Big(\sum_{i,j}\ketbra ij\otimes\ketbra{\tilde\imath}{\tilde\jmath}\Big)=\sum_{i,j}\braket{\tilde\jmath}{\tilde\imath}\ketbra ij,
\end{equation}
because, using the completeness relation $\sum_k\proj k=I$ in $\Ha_B$, $\tr_B\ketbra{\tilde\imath}{\tilde\jmath}=\sum_k\braket k{\tilde\imath}\braket{\tilde\jmath}k=\braket{\tilde\jmath}{\tilde\imath}$. Comparing the two expressions of $\hat\rho_A$,
\begin{equation}
  \braket{\tilde\jmath}{\tilde\imath}_B=p_i\delta_{ij}:
\end{equation}
the states $\ket{\tilde\imath}_B$ are orthogonal. We can therefore normalize them, $\ket{i'}_B=p_i^{-1/2}\ket{\tilde\imath}_B$ (for $p_i\neq0$), and obtain
\begin{equation}\label{eq:schmidt}
  \ket\psi_{AB}=\sum_i\sqrt{p_i}\ket i_A\otimes\ket{i'}_B .
\end{equation}
This is the \emph{Schmidt decomposition}: any pure state of a bipartite system can be written as a single sum over pairs of orthonormal states $\ket i_A$, $\ket{i'}_B$ of the two subsystems, with non-negative coefficients $\sqrt{p_i}$, the
\emph{Schmidt coefficients}. The bases depend on the state. Compared with the general expansion, which has $d_Ad_B$ terms, the Schmidt form has at most $\min(d_A,d_B)$ terms.

\paragraph{Singular value decomposition.} In fixed bases $\{\ket a_A\}$ and $\{\ket\mu_B\}$, the Schmidt decomposition reads
\begin{equation}
  \psi_{a\mu}=\sum_i(U_A)_{ai}\sqrt{p_i}\,(U_B^T)_{i\mu},
\end{equation}
where the unitary matrices $U_A$ and $U_B$ transform the fixed bases into the Schmidt bases. This is the singular value decomposition (SVD) $\Psi=U_A\Sigma U_B^T$ of the matrix $\Psi=(\psi_{a\mu})$ of coefficients. Conversely, the SVD theorem
(any complex matrix is $U\Sigma V^\dagger$ with $U,V$ unitary and $\Sigma\ge0$ diagonal) gives a second proof of the Schmidt decomposition. The Schmidt coefficients are the singular values of $\Psi$, which gives a practical way to compute them
numerically (\texttt{numpy.linalg.svd}). The reduced density matrices are $\hat\rho_A=\Psi\Psi^\dagger$ and $\hat\rho_B=(\Psi^\dagger\Psi)^T$ in these bases.

\section{Consequences}
\paragraph{Equal spectra.} Tracing out $A$ instead of $B$ gives
\begin{equation}
  \hat\rho_B=\tr_A\proj\psi=\sum_ip_i\proj{i'}.
\end{equation}
The reduced density matrices $\hat\rho_A$ and $\hat\rho_B$ therefore have the same non-zero eigenvalues $p_i$, even if the two subsystems have different dimensions. In particular they have the same purity and the same von Neumann entropy,
which justifies the definition of the entanglement entropy $S(\hat\rho_A)=S(\hat\rho_B)=-\sum_ip_i\log_2p_i$.

\paragraph{Schmidt number and entanglement.} The number of non-zero Schmidt coefficients is the \emph{Schmidt number} (the rank of $\Psi$). A state is entangled if and only if its Schmidt number is greater than 1. For a separable state,
$\ket\psi_{AB}=\ket\varphi_A\otimes\ket\chi_B$, there is a single term, and the reduced density matrices $\hat\rho_A=\proj\varphi$, $\hat\rho_B=\proj\chi$ describe pure states. Otherwise, they are mixed, and subsystem $A$ alone has no wave function.
For example, the Bell state $(\ket{00}+\ket{11})/\sqrt2$ is already in Schmidt form, with two coefficients $1/\sqrt2$: its Schmidt number is 2, and $\hat\rho_A=\hat\rho_B=I/2$, in agreement with the partial trace computed in \cref{sec:reduced}.
For two qubits, $\det\Psi=\sqrt{p_1p_2}\det U_A\det U_B$, so $|\det\Psi|=\sqrt{p_1p_2}$, which explains the criterion of \cref{prop:det} and the concurrence $C=2\sqrt{p_1p_2}$.

\paragraph{Entanglement cannot be created locally.} A unitary operation $U_A\otimes U_B$ acting separately on the two subsystems only rotates the Schmidt bases,
\begin{equation}
  (U_A\otimes U_B)\ket\psi_{AB}=\sum_i\sqrt{p_i}\left(U_A\ket i_A\right)\otimes\left(U_B\ket{i'}_B\right),
\end{equation}
and leaves the Schmidt coefficients unchanged. In particular, local operations cannot increase the Schmidt number. Entanglement can only be created by an interaction between the subsystems, such as a CNOT gate. (More generally, entanglement cannot
increase on average under local operations and classical communication; Nielsen's theorem (1999) states that $\ket\psi$ can be converted into $\ket\phi$ by LOCC iff the Schmidt coefficients of $\ket\psi$ are majorized by those of $\ket\phi$.)

\section{Purification}\label{app:purification}
Conversely, any mixed state can be seen as part of a larger pure state. Suppose that $\hat\rho_A$ does not describe a pure state. We can introduce a fictitious \emph{reference system} $R$ such that the joint state of $A$ and $R$ is pure and
\begin{equation}
  \hat\rho_A=\tr_R\proj{AR}.
\end{equation}
\begin{proof}
Write $\hat\rho_A$ in its orthonormal eigenbasis, $\hat\rho_A=\sum_ip_i\proj{i^A}$. Let $R$ have the same state space as $A$, with an orthonormal basis $\{\ket{i^R}\}$, and define
\begin{equation}
  \ket{AR}=\sum_i\sqrt{p_i}\ket{i^A}\ket{i^R}.
\end{equation}
Then
\begin{equation}
  \tr_R\proj{AR}=\sum_{i,j}\sqrt{p_ip_j}\ketbra{i^A}{j^A}\tr\left(\ketbra{i^R}{j^R}\right)=\sum_{i,j}\sqrt{p_ip_j}\ketbra{i^A}{j^A}\delta_{ij}=\sum_ip_i\proj{i^A}=\hat\rho_A.\qedhere
\end{equation}
\end{proof}
The state $\ket{AR}$ is written directly in Schmidt form, which shows the close relation between purification and the Schmidt decomposition. Physically, purification says that mixedness can always be attributed to entanglement with
something else (the ``church of the larger Hilbert space''). A compact formula: with the unnormalized maximally entangled vector $\ket\Omega=\sum_i\ket i\ket i$, every purification of $\hat\rho$ is $(\sqrt{\hat\rho}\otimes V)\ket\Omega$
for some unitary (or isometry) $V$ on $R$. This is used in the proof of Uhlmann's theorem (\cref{app:fidelity}).

\section{The Hughston--Jozsa--Wootters theorem}\label{app:hjw}
A given density matrix can be realized by many different ensembles of pure states \cite{HJW1993}:
\begin{equation}
  \hat\rho_A=\sum_ip_i\proj{\varphi_i}=\sum_\mu q_\mu\proj{\psi_\mu},
\end{equation}
where neither the $\ket{\varphi_i}$ nor the $\ket{\psi_\mu}$ need to be orthogonal. Each ensemble corresponds to a purification,
\begin{equation}
  \ket{\Phi_1}_{AB}=\sum_i\sqrt{p_i}\ket{\varphi_i}_A\ket{\alpha_i}_B,\qquad\ket{\Phi_2}_{AB}=\sum_\mu\sqrt{q_\mu}\ket{\psi_\mu}_A\ket{\beta_\mu}_B,
\end{equation}
where $\{\ket{\alpha_i}\}$ and $\{\ket{\beta_\mu}\}$ are orthonormal states of a system $B$ of sufficiently large dimension. (Check: $\tr_B\proj{\Phi_1}=\sum_{ij}\sqrt{p_ip_j}\ketbra{\varphi_i}{\varphi_j}\braket{\alpha_j}{\alpha_i}=\hat\rho_A$.)
\begin{theorem}[HJW]
Two purifications of the same $\hat\rho_A$ on the same space $B$ are related by a unitary change of basis in $B$ alone:
\begin{equation}
  \ket{\Phi_1}_{AB}=(I_A\otimes U_B)\ket{\Phi_2}_{AB}.
\end{equation}
\end{theorem}
\begin{proof}
Both states give the same reduced density matrix $\hat\rho_A$, so their Schmidt decompositions can be written with the same coefficients and the same eigenbasis $\{\ket k_A\}$ of $\hat\rho_A$ (if some eigenvalues are degenerate, we choose the same basis
of each eigenspace for both; this is possible because \cref{eq:schmidt} holds for any eigenbasis of $\hat\rho_A$):
\begin{equation}
  \ket{\Phi_1}_{AB}=\sum_k\sqrt{\lambda_k}\ket k_A\ket{k_1}_B,\qquad\ket{\Phi_2}_{AB}=\sum_k\sqrt{\lambda_k}\ket k_A\ket{k_2}_B .
\end{equation}
The linear map defined by $\ket{k_1}_B=U_B\ket{k_2}_B$ maps one orthonormal set onto another; it can be extended to a unitary on the whole of $\Ha_B$, which maps one state onto the other.
\end{proof}
An equivalent algebraic form (\cite{NielsenChuang}, Thm.~2.6): $\sum_ip_i\proj{\varphi_i}=\sum_\mu q_\mu\proj{\psi_\mu}$ iff $\sqrt{p_i}\ket{\varphi_i}=\sum_\mu u_{i\mu}\sqrt{q_\mu}\ket{\psi_\mu}$ for some unitary matrix $u$ (padding the
smaller ensemble with zero-probability states).

The physical meaning is remarkable. If $A$ and $B$ share the state $\ket{\Phi_1}$, then by choosing in which basis to measure $B$, an observer can decide which ensemble describes $A$: measuring $B$ in the basis $\{\ket{\alpha_i}\}$ leaves $A$ in
$\ket{\varphi_i}$ with probability $p_i$; measuring in the basis $\{U_B\ket{\beta_\mu}\}$ leaves $A$ in $\ket{\psi_\mu}$ with probability $q_\mu$. Since both ensembles give the same $\hat\rho_A$, an observer of $A$ alone cannot tell which choice was made,
and no information is transmitted. This is the precise content of ``steering'' (Schrödinger, 1935) and of the EPR argument (\cref{ch:bell}), and again of the no-signalling principle.

\section{Exercises}
\begin{exercise}
Find the Schmidt decomposition of $\frac12(\ket{00}+\ket{01}+\ket{10}-\ket{11})$. (Hint: compute the SVD of $\frac12\begin{pmatrix}1&1\\1&-1\end{pmatrix}$.) Is the state entangled?
\end{exercise}
\begin{exercise}
Show that the two ensembles $\{\frac12,\ket0;\frac12,\ket1\}$ and $\{\frac12,\ket+;\frac12,\ket-\}$ are related by the unitary $u=H$ in the HJW theorem, and find the corresponding measurements on $B$ for the purification $\ket{\Phi^+}$.
\end{exercise}
